\documentclass[11pt]{article}
\usepackage[a4paper,margin=22mm]{geometry}
\usepackage{amsmath,amssymb,amsthm,mathtools,bm}
\usepackage{booktabs,longtable,array}
\usepackage{enumitem}
\usepackage{xcolor,microtype,xurl}
\usepackage[hidelinks]{hyperref}
\setlength{\emergencystretch}{3em}
\numberwithin{equation}{section}
\newtheorem{theorem}{Theorem}[section]
\newtheorem{proposition}[theorem]{Proposition}
\newtheorem{lemma}[theorem]{Lemma}
\newtheorem{corollary}[theorem]{Corollary}
\theoremstyle{definition}
\newtheorem{definition}[theorem]{Definition}
\newtheorem{gate}[theorem]{Fail-closed gate}
\theoremstyle{remark}
\newtheorem{remark}[theorem]{Remark}
\newcommand{\RR}{\mathbb R}
\newcommand{\TT}{\mathbb T}
\newcommand{\SSS}{\mathbb S}
\newcommand{\cB}{\mathcal B}
\newcommand{\cE}{\mathcal E}
\newcommand{\cT}{\mathcal T}
\newcommand{\dd}{\mathrm d}
\newcommand{\ezer}{e_0}
\newcommand{\whom}{W_{\mathrm{hom}}}

\title{AP-E10 Exact Stencil No-Go and Finite-$R$ Cell Formula:\\
Discrete Skyrme Defects, Periodic Homogenization,\\
and a Translation-Quotient Continuum Audit}
\author{Route F research note}
\date{18 July 2026}

\begin{document}
\maketitle

\begin{abstract}
This note closes two AP-E9 questions and gives a negative answer to the first
physical conjunction.  The one-corner forward Skyrme stencil does not satisfy
the compensated-compactness property required by the proposed continuum
argument.  An exact three-periodic sequence $n_a:\TT^3\to\SSS^3$ converges
uniformly to the vacuum while one of its forward two-minors remains the
nonzero constant $3\sqrt3\eta^2/2$.  A second two-periodic sequence has
geometric degree zero but forward topological current converging to
$-16\eta^3$.  Localizing it with an exact vacuum boundary produces the diffuse
defect $-16\eta^3\chi^3\dd x$.  On the other hand the Skyrme bound implies
$|J_a|^{4/3}\le |\wedge^2D^an_a|^2/3$, so singular concentration of this
forward current is excluded.  The obstruction is a diffuse lattice-scale
aliasing defect, not loss of mass into an atomic measure.

The finite-$R$ barrier problem is solved exactly.  On both the uniform
six-tetrahedron graph and the two checkerboard five-tetrahedron phases, every
bond-direction subgraph is cycle-balanced.  Direction-wise Jensen inequality
therefore proves that the periodic corrector is zero and that the AP-E9 raw
Cauchy--Born quartics are the homogenized densities.  Their ratios are $4/3$
and $3$ on two gradients, so triangulation dependence survives homogenization.
An eleven-case production scan then adds a physical translation quotient,
dynamic centering, $a\downarrow0.20833$, and $L\uparrow7.5$.  All cases are
stationary, admissible, and have three-target degree one, but size, radial
profile, and cross-triangulation gates fail.  Consequently the complete
Gamma-liminf, continuum background, same-action Hessian, and regulated
determinant variation remain prohibited.  The exact repair is structural: a
compatible discrete de Rham/cup-product Skyrme term must replace the
one-corner product before the continuum program can resume.
\end{abstract}

\tableofcontents

\section{Declared stencil and proof obligations}

Let $a>0$ and write the forward differences on a cubic lattice as
\begin{equation}
 D_i^a n(z)=\frac{n(z+e_i)-n(z)}{a},\qquad i=1,2,3.
 \label{eq:forward}
\end{equation}
The AP-E7 Skyrme density at the lower corner of a cube is
\begin{equation}
 S_a(z)=\sum_{1\le i<j\le3}
 \left|D_i^an(z)\wedge D_j^an(z)\right|^2.
 \label{eq:skyrme-density}
\end{equation}
It is sampled once per cube rather than assembled from the gradients of all
tetrahedral interpolants.  The associated algebraic forward current is
\begin{equation}
 J_a(z)=\det\bigl[n(z),D_1^an(z),D_2^an(z),D_3^an(z)\bigr].
 \label{eq:forward-current}
\end{equation}

For genuine gradients, minors are null Lagrangians and div--curl arguments can
give weak continuity under appropriate hypotheses
\cite{BrianeCasadoDiaz2016}.  Such an argument is not automatic for a finite
difference product: the ordinary product rule contains shifts, and the
unshifted expression in Eq.~\eqref{eq:skyrme-density} is not a cochain cup
product.  Compatible discretizations instead preserve a differential complex
and a bounded cochain projection \cite{ArnoldFalkWinther2010}.  AP-E10 tests
the actual AP-E7 stencil, not an ideal compatible replacement.

Three statements must be separated:
\begin{enumerate}[label=(\roman*)]
\item weak continuity of the two-minors in Eq.~\eqref{eq:skyrme-density};
\item exclusion of a singular concentration measure for $J_a$;
\item identification of the weak limit of $J_a$ with the geometric degree
current of the normalized-affine map.
\end{enumerate}
The first and third statements fail; a quantitative version of the second is
true.

\section{Exact failure of compensated compactness}

\subsection{A three-periodic two-minor defect}

Work first on the unit torus with $a=(3m)^{-1}$.  For lattice index
$z=(z_1,z_2,z_3)$ put
\begin{equation}
 \theta_z=\frac{2\pi}{3}(z_1+2z_2),\qquad
 n_a(z)=\left(\sqrt{1-\eta^2a^2},
 \eta a\cos\theta_z,\eta a\sin\theta_z,0\right),
 \label{eq:minor-sequence}
\end{equation}
where $\eta>0$ is fixed and $a<\eta^{-1}$.  The sequence is exactly
sphere-valued and $n_a\to\ezer$ uniformly.

\begin{lemma}[Forward phase identity]
Let $f_z=\cos\theta_z$, $g_z=\sin\theta_z$, and suppose a forward step in
direction $i$ changes the phase by $\alpha_i$.  Then
\begin{equation}
 (\Delta_i f)(\Delta_j g)-(\Delta_i g)(\Delta_j f)
 =4\sin\frac{\alpha_i}{2}\sin\frac{\alpha_j}{2}
   \sin\frac{\alpha_j-\alpha_i}{2}.
 \label{eq:phase-identity}
\end{equation}
\end{lemma}

\begin{proof}
Use
$\Delta_i f=-2\sin(\alpha_i/2)\sin(\theta+\alpha_i/2)$ and
$\Delta_i g=2\sin(\alpha_i/2)\cos(\theta+\alpha_i/2)$, expand the determinant,
and apply the sine subtraction identity.  The result is independent of
$\theta$.
\end{proof}

Here $\alpha_1=2\pi/3$ and $\alpha_2=4\pi/3$.  Therefore the target-component
$(1,2)$ entry of the spatial $(1,2)$ Skyrme minor is exactly
\begin{equation}
 \left(D_1^an\wedge D_2^an\right)_{12}
 =\frac{3\sqrt3}{2}\eta^2
 \label{eq:minor-constant}
\end{equation}
at every lattice site.

\begin{theorem}[One-corner compensated compactness is false]
There exists an $\SSS^3$-valued sequence converging uniformly, hence strongly
in every finite $L^p$, to the vacuum, whose AP-E7 forward two-minors converge
weakly to a nonzero constant although the continuum minor of the limit is
zero.  Its discrete Dirichlet and Skyrme energies remain uniformly bounded.
\end{theorem}

\begin{proof}
The sequence \eqref{eq:minor-sequence} has
$\|n_a-\ezer\|_{L^\infty}=O(a)$.  Its forward derivatives are $O(\eta)$ and
its two-minors are $O(\eta^2)$; summation with cell volume $a^3$ therefore
gives uniform Dirichlet and Skyrme bounds.  Equation
\eqref{eq:minor-constant} contradicts weak continuity.  Multiplication of the
tangent components by a fixed $\chi\in C_c^\infty(\Omega)$ followed by exact
normalization gives an exact-vacuum-boundary version whose defect tends to
$(3\sqrt3/2)\eta^2\chi^2$ on the support of $\chi$.
\end{proof}

For $\eta=0.2$, the verifier obtains the same value
$0.1039230484541327$ at $a=1/6,1/12,1/24,1/48$, with maximum exact residual
$5.6\times10^{-17}$ while the uniform distance to the vacuum decreases to
$4.17\times10^{-3}$.

\subsection{A two-periodic diffuse topological-current defect}

On an even lattice set
\begin{equation}
 u_1=(-1)^{z_1+z_2},\qquad
 u_2=(-1)^{z_2+z_3},\qquad
 u_3=(-1)^{z_3+z_1}
 \label{eq:checker-u}
\end{equation}
and define
\begin{equation}
 n_a(z)=\left(\sqrt{1-3\eta^2a^2},\eta a u_1,\eta a u_2,\eta a u_3\right).
 \label{eq:current-sequence}
\end{equation}
Again $n_a\to\ezer$ uniformly and the image lies in the open northern
hemisphere, so every geometric degree is zero.  The tangent derivative matrix
is
\begin{equation}
 \eta\begin{pmatrix}
 -2u_1&-2u_1&0\\
 0&-2u_2&-2u_2\\
 -2u_3&0&-2u_3
 \end{pmatrix}.
 \label{eq:checker-matrix}
\end{equation}
Since $u_1u_2u_3=1$, its determinant is $-16\eta^3$.  Thus
\begin{equation}
 \boxed{J_a=-16\eta^3\sqrt{1-3\eta^2a^2}
 \longrightarrow -16\eta^3,}
 \label{eq:diffuse-current}
\end{equation}
while the geometric degree and continuum current of the limit are zero.

For a fixed-boundary form, let
\begin{equation}
 \chi(x)=\prod_{i=1}^3\sin^2(\pi x_i)
 \label{eq:cutoff}
\end{equation}
on $[0,1]^3$, replace the tangent part by $\eta a\chi u$, and normalize.
Then the boundary is exactly the vacuum and
\begin{equation}
 J_a\rightharpoonup -16\eta^3\chi^3\dd x,\qquad
 \int J_a\dd x\longrightarrow
 -16\eta^3\left(\frac5{16}\right)^3.
 \label{eq:localized-current}
\end{equation}
For $\eta=0.2$ the limit is $-0.00390625$.  At $a=1/8$ and $1/48$ the
computed integrals are respectively $-0.00295055$ and $-0.00387605$; the
Dirichlet and Skyrme energies stay below $0.026$ and $0.0024$.

\section{What the Skyrme bound still proves about concentration}

Let $A_a$ be the $4\times3$ matrix with columns $D_i^an$, and let
$s_1,s_2,s_3$ be its singular values.  Pointwise,
\begin{align}
 |J_a|&\le |\wedge^3A_a|=s_1s_2s_3,\\
 |\wedge^2A_a|^2
 &=s_1^2s_2^2+s_1^2s_3^2+s_2^2s_3^2.
\end{align}
Arithmetic--geometric mean gives
\begin{equation}
 |\wedge^2A_a|^2\ge
 3(s_1s_2s_3)^{4/3}.
 \label{eq:amgm}
\end{equation}

\begin{theorem}[No singular forward-current concentration]
Every sequence with uniformly bounded AP-E7 Skyrme energy satisfies
\begin{equation}
 \boxed{\int_\Omega |J_a|^{4/3}\dd x
 \le\frac13\int_\Omega|\wedge^2D^an_a|^2\dd x.}
 \label{eq:l43}
\end{equation}
Hence $J_a$ is weakly precompact in $L^{4/3}$ and cannot acquire an atomic or
other singular-measure concentration defect.
\end{theorem}

\begin{remark}
Equation \eqref{eq:l43} excludes concentration but does not identify the weak
limit.  Equations \eqref{eq:minor-constant} and \eqref{eq:diffuse-current}
show that a diffuse oscillatory defect survives.  Nor does
Eq.~\eqref{eq:l43} control the geometric normalized-affine current, because
the AP-E7 action samples only one forward triad per cube.  A theorem for
geometric degree requires a compatible simplex-wise action and current.
Recent two-dimensional charge-measure compactness results
\cite{BrianiCicaleseKreutz2026} demonstrate the needed architecture but do
not transfer to this three-dimensional, noncompatible stencil.
\end{remark}

\section{Finite-R periodic cell problem}

At the AP-E9 critical scale let $R_a=\gamma_a a^2/d_a^2\to R\in(0,\infty)$.
Around a fixed target point, a bounded microscopic tangent corrector gives a
bond increment $a(Ar+D_r\varphi)+o(a)$.  The logarithmic barrier expansion
therefore yields
\begin{equation}
 \cB_a\longrightarrow \frac{R}{8}\int_\Omega
 Q_{\mathrm{hom}}(\nabla n)\dd x.
 \label{eq:barrier-limit}
\end{equation}

Let $Y=(\mathbb Z/P\mathbb Z)^3$ and let $\mathcal B_Y$ be the unique periodic
bonds per cell, written as $(p,r)$ from $p$ to $p+r$.  The exact convex cell
problem is
\begin{equation}
 Q_{\mathrm{hom}}(A)=\frac1{|Y|}
 \inf_{\varphi:Y\to\RR^m}
 \sum_{(p,r)\in\mathcal B_Y}
 |Ar+\varphi(p+r)-\varphi(p)|^4.
 \label{eq:cell-problem}
\end{equation}
The additive constant of $\varphi$ is immaterial.

\begin{definition}[Direction-wise cycle balance]
For each bond direction $r$, let $\mathcal B_{Y,r}$ be its directed periodic
subgraph.  It is cycle-balanced if every vertex has the same incoming and
outgoing multiplicity.
\end{definition}

\begin{lemma}[Zero-corrector criterion]
If every $\mathcal B_{Y,r}$ is cycle-balanced, then the zero corrector solves
Eq.~\eqref{eq:cell-problem} exactly.
\end{lemma}

\begin{proof}
Cycle balance gives
\begin{equation}
 \sum_{(p,r)\in\mathcal B_{Y,r}}
 [\varphi(p+r)-\varphi(p)]=0.
\end{equation}
Because $v\mapsto|v|^4$ is convex, Jensen's inequality gives
\begin{equation}
 \frac1{|\mathcal B_{Y,r}|}\sum_{(p,r)\in\mathcal B_{Y,r}}
 |Ar+D_r\varphi(p)|^4\ge |Ar|^4.
\end{equation}
Summing directions proves the lower bound, and $\varphi=0$ attains it.
\end{proof}

The uniform six-tet graph has, per primitive cell, the seven directions
\begin{equation}
 e_1,e_2,e_3,e_1+e_2,e_1+e_3,e_2+e_3,e_1+e_2+e_3.
\end{equation}
The checkerboard five-tet graph has the three coordinate directions and both
signs of every face diagonal with weight $1/2$.  Exact enumeration on the
period-two cell gives nodewise in-degree equal to out-degree for every one of
these direction graphs.  Thus
\begin{align}
 Q_6^{\mathrm{hom}}(A)
 &=\sum_i|Ae_i|^4+\sum_{i<j}|A(e_i+e_j)|^4
   +|A(e_1+e_2+e_3)|^4,
 \label{eq:q6}\\
 Q_5^{\mathrm{hom}}(A)
 &=\sum_i|Ae_i|^4+\frac12\sum_{i<j}
 \bigl(|A(e_i+e_j)|^4+|A(e_i-e_j)|^4\bigr).
 \label{eq:q5}
\end{align}
The two checkerboard phases are lattice translates and have the same
$Q_5^{\mathrm{hom}}$.  Random-start numerical corrector solves for four fixed
and four random $3\times3$ gradients reproduce Eqs.~\eqref{eq:q6}--
\eqref{eq:q5} with maximum relative residual $9.1\times10^{-15}$.

\begin{theorem}[Finite-$R$ mesh dependence survives homogenization]
The six- and five-tet homogenized densities are not scalar multiples.  If
$(Ae_1,Ae_2,Ae_3)=(v,0,0)$, then
$Q_6^{\mathrm{hom}}/Q_5^{\mathrm{hom}}=4/3$; if it is $(v,v,0)$, the ratio is
$3$.  No renormalization of the single coefficient $R$ can equate them.
\end{theorem}

This closes the AP-E9 cell problem, but negatively: a finite critical barrier
is a mesh-dependent physical quartic.  General discrete homogenization
frameworks \cite{AlicandroBraidesCicaleseSolci2023} justify the cell-formula
architecture; here the cycle certificate makes the minimization explicit.

\section{Translation quotient and dynamic centering}

For a localized soliton define the nonnegative density, barycentre, and
covariance
\begin{equation}
 \rho=1-n_0,\qquad
 c[n]=\frac{\int x\rho\dd x}{\int\rho\dd x},\qquad
 \Sigma[n]=\frac{\int(x-c)(x-c)^T\rho\dd x}{\int\rho\dd x}.
 \label{eq:quotient}
\end{equation}
The translation quotient compares $\operatorname{tr}\Sigma$, its eigenvalues,
and the normalized radial histogram about $c[n]$.  It therefore identifies
fields differing only by physical translation without adding an optimizer
anchor.  This distinction matters on growing boxes, for which ordinary
strong-$L^2$ equicoercivity fails unless translations are quotiented or made
tight.

The production driver also attempts a dynamic representative: after a
stationary solve it shifts the field by $-c[n]$ using normalized linear
interpolation, accepts the shift only if the logarithmic edge condition and
all three geometric degree targets remain valid, and then re-relaxes the
same action.  Re-relaxation can return the field to a mesh-preferred Peierls
position.  Consequently the physical quotient, not a small raw barycentre,
is the primary translation gate.  Convergence of harmonic-map discretizations
likewise requires declared triangulation and weight hypotheses rather than
an informal refinement argument \cite{GasterLoustauMonsaingeon2023}.

\section{Production scan}

The scan contains eleven independent unanchored solutions.  The joint
sequence is
\begin{center}
\begin{tabular}{rrrrrr}
\toprule
$N$&$L$&$a$&total energy&centered RMS&gradient density\\
\midrule
25&6.0&0.250000&76.672058&1.18660&$3.04\times10^{-7}$\\
29&6.5&0.232143&77.679545&1.22405&$1.64\times10^{-6}$\\
33&7.0&0.218750&78.552231&1.26582&$5.34\times10^{-7}$\\
37&7.5&0.208333&78.987847&1.35422&$1.70\times10^{-6}$\\
\bottomrule
\end{tabular}
\end{center}
All eleven cases have direct tangent-plus-scalar gradient density below
$2\times10^{-6}$, positive edge margin at least $0.1064$, and three-target
$B=(1,1,1)$.  Production coverage and mechanical checks pass.  The
predeclared quotient diagnostics are
\begin{center}
\begin{tabular}{lrrr}
\toprule
diagnostic&observed&threshold&pass\\
\midrule
joint-tail energy spread&1.656\%&3\%&yes\\
joint-tail centered-RMS spread&9.612\%&5\%&no\\
joint maximum successive profile $L^1$&7.072\%&5\%&no\\
translated-start energy spread&0.00192\%&1\%&yes\\
translated-start centered-RMS spread&0.0602\%&2\%&yes\\
translated-start profile $L^1$&3.500\%&3\%&no\\
five-/six-tet energy spread&5.143\%&3\%&no\\
five-/six-tet centered-RMS spread&11.378\%&5\%&no\\
\bottomrule
\end{tabular}
\end{center}
The maximum post-relaxation barycentre residual is $0.740a$, confirming
Peierls locking and the need for the quotient.  The cross-mesh discrepancy
agrees in sign with the exact finite-$R$ no-go and is not an optimizer
failure.

\section{Fail-closed decision and structural repair}

\begin{gate}[Continuum and Hessian authorization]
The current AP-E7/AP-E8 lattice action cannot authorize a full Gamma-liminf,
regulator-stable continuum background, Riemann Hessian, or regulated
determinant variation because:
\begin{enumerate}
\item the one-corner two-minors are not weakly continuous;
\item the forward topological current has a nonzero diffuse defect at
geometric degree zero;
\item finite-$R$ five- and six-tet homogenized quartics are inequivalent;
\item the joint size/profile and cross-mesh numerical gates fail.
\end{enumerate}
\end{gate}

Thus
\begin{align}
 \texttt{forward\_minor\_compensated\_compactness}&=\texttt{false},\\
 \texttt{finite\_R\_cell\_density\_computed}&=\texttt{true},\\
 \texttt{finite\_R\_triangulation\_independence}&=\texttt{false},\\
 \texttt{regulator\_stable\_background}&=\texttt{false},\\
 \texttt{same\_action\_Hessian\_allowed}&=\texttt{false}.
\end{align}

The next action should change the stencil rather than extend the same scan:
\begin{enumerate}
\item define $dn$ as a primal one-cochain on the complete tetrahedral complex;
\item form the Skyrme two-form with a graded, shifted cup product satisfying a
discrete Leibniz rule and $d^2=0$;
\item use a positive discrete Hodge star, averaged over the four body-diagonal
six-tet meshes or calibrated by a cubic fourth-order design;
\item define the topological three-cochain from the same cup product and prove
that its evaluation equals normalized-affine degree on admissible fields;
\item rerun the cell certificate and only then repeat the quotient continuum
scan.
\end{enumerate}
The FEEC principle is that stability requires a subcomplex and bounded
cochain projection \cite{ArnoldFalkWinther2010}; this is the appropriate
first-principles repair.  Historical lattice-Skyrme work already warns that a
topology-preserving cutoff can become part of the model
\cite{SchrammSvetitsky2000}.  The finite GW/domain-wall kernel and actual
$SO(3)$ mapping-torus mod-two index remain parallel lanes, and the degree-one
Route-E portal remains last.

\section{Reproducibility}

Run from the repository root:
\begin{verbatim}
PYTHONPYCACHEPREFIX=/tmp/ap_e10_pycache python3 \
  route_f/code/scan_ap_e10_compactness_homogenization_centering.py
\end{verbatim}
The deterministic result is written to
\begin{center}
\path{route_f/output/ap_e10_compactness_homogenization_centering.json}
\end{center}
and its Markdown companion.  The production card passes $10/10$ exact,
numerical, provenance, and fail-closed checks.  A pass means that each claimed
calculation was executed; it does not turn a deliberately false physics gate
true.

\bibliographystyle{unsrt}
\bibliography{route_f/tex/ap_e10_compactness_homogenization_centering}

\end{document}
